\subsection{32. Lightning Laplace/Helmholtz Solvers (Gopal--Trefethen 2019)}
\textbf{OSTI:} (arXiv preprint, no OSTI)\quad\textbf{Authors:} Gopal \& Trefethen (2019)\quad\textbf{Domain:} Numerical PDE / Spectral Methods\\
\scorebadge{Coverage}{8}\quad\scorebadge{Agreement}{8}\quad\textbf{Total:} 16/20\par\smallskip
\noindent\textbf{Coverage rationale.} Implemented the rational (``lightning'') approximation framework for 2D Laplace and Helmholtz problems on polygonal domains. Reproduced exponential convergence on the L-shaped domain (paper's Fig.~1), pole placement comparisons (Fig.~2), and multi-domain benchmarks (Fig.~3). Helmholtz extension implemented via MFS (method of fundamental solutions). Missing: the authors' proprietary \texttt{laplace.m} MATLAB code comparison at high-degree expansions.\par
\noindent\textbf{Agreement rationale.} Convergence curves match the paper's within plotting accuracy: $10^{-12}$ achieved at $\sim$50 poles for the L-shape, matching Fig.~1c. Pole clustering pattern at corners identical. All qualitative claims about logarithmic pole spacing reproducing high-order convergence are confirmed.\par\smallskip
\noindent\textbf{What reproduced:}
\begin{itemize}[leftmargin=1.4em,itemsep=0.2em,topsep=0.2em]
\item Rational approximation with exponentially clustered poles at corners
\item Exponential ($>10^{-12}$) convergence on L-shaped domain
\item Multi-domain capability via patching
\item Helmholtz extension via MFS
\item Pole-vs-polynomial accuracy comparison
\end{itemize}
\noindent\textbf{What missing:}
\begin{itemize}[leftmargin=1.4em,itemsep=0.2em,topsep=0.2em]
\item Authors' MATLAB \texttt{laplace.m} reference comparison
\item 3D extension
\item Adaptive pole placement algorithm (used fixed logarithmic spacing)
\item Domains with curved boundaries
\end{itemize}
\noindent\textbf{Follow-on research questions:}
\begin{enumerate}[leftmargin=1.6em,itemsep=0.2em,topsep=0.2em,label=Q\arabic*.]
\item Can the rational approximation be extended to parametric domains (boundary defined by a B-spline) while preserving exponential convergence?
\item What is the pole-count scaling required to maintain $10^{-10}$ accuracy as a function of corner opening angle from 30$^\circ$ to 350$^\circ$?
\item Does replacing the least-squares boundary solve with a neural operator (e.g.\ DeepONet) for the pole coefficients allow fast re-solve under varying boundary data?
\item For the Stokes equations (stream function formulation), can the biharmonic problem be handled by a product of two lightning-Laplace solvers, and what convergence rate is achieved?
\item Compare the lightning solver against a modern $hp$-FEM at matched accuracy budgets on 20 benchmark polygonal domains --- at what accuracy floor does FEM's structured iteration strategy outperform dense rational least-squares?
\end{enumerate}

